\documentclass[12pt,man]{apa7}
\usepackage{styles/cwilliams-apa}
\usepackage{booktabs}
\usepackage{graphicx}
\usepackage{multirow}
\usepackage{amsmath}
\usepackage{float}

\title{Machine Learning Term Project Report: Predicting Traffic Accident Injury Severity Using Comprehensive Machine Learning Techniques}
\shorttitle{Traffic Accident ML Analysis}

\author{Christopher Williams}
\affiliation{Virginia Tech}

\course{5805: Machine Learning I}
\professor{Dr.\ Reza Jafari}
\duedate{December 6, 2025}

\begin{document}

\maketitle

\tableofcontents

\listoffigures
\listoftables

\section{Abstract}

Traffic accidents represent one of the most significant contributors to injuries, fatalities, and economic losses worldwide. This study applies comprehensive machine learning techniques to predict crash injury severity using 209,275 traffic accident records spanning 12 years (2013-2025). The analysis follows a systematic four-phase approach encompassing feature engineering, regression analysis, classification, and unsupervised learning methods. Multiple linear regression on the continuous target variable \texttt{injuries\_total} achieved an $R^2$ of 0.066 with 75\% of coefficients statistically significant (p $<$ 0.05). For binary classification of \texttt{crash\_type}, Random Forest with bagging demonstrated superior performance with 74.06\% accuracy, 74.03\% F1-score, and 0.81 AUC, outperforming nine other classifiers including LDA, Decision Trees, Logistic Regression, KNN, SVM, Naive Bayes, ensemble methods, and Neural Networks. Feature importance analysis revealed that \texttt{first\_crash\_type}, \texttt{damage}, \texttt{prim\_contributory\_cause}, \texttt{num\_units}, and \texttt{failure\_to\_obey} were the most significant predictors. Clustering analysis using K-Means++ identified K=2 optimal clusters with silhouette score of 0.3481, while association rule mining discovered 1,037 rules with the strongest rule (RAIN $\rightarrow$ WET road surface, lift=6.16) demonstrating high predictive relationships. The findings provide actionable insights for traffic safety policies, emergency resource allocation, and accident prevention strategies. However, limited data granularity (binary target, coarse temporal features, lack of spatial coordinates) constrains model performance, suggesting future work should incorporate finer-grained variables such as GPS coordinates, vehicle speeds, driver demographics, and traffic volume data.

\section{Introduction}

\subsection{Background and Motivation}

Traffic accidents pose a critical public safety challenge, resulting in approximately 1.35 million deaths globally each year according to the World Health Organization. In the United States alone, traffic crashes cost over \$340 billion annually in medical expenses, property damage, and lost productivity. Understanding the factors that contribute to accident severity is essential for developing effective prevention strategies, optimizing emergency response protocols, and informing transportation infrastructure design.

Despite advances in vehicle safety technology and traffic management systems, predicting crash outcomes remains challenging due to the complex interplay of environmental conditions, road characteristics, driver behavior, and vehicle factors. Machine learning offers powerful tools to uncover hidden patterns in large-scale accident datasets and build predictive models that can assist traffic safety officials and emergency services.

\subsection{Project Objectives}

The primary goal of this Final Term Project (FTP) is to obtain practical experience with feature engineering, classification, regression, clustering, and association rule mining by applying various machine learning algorithms to real-world traffic accident data. Specifically, this analysis aims to:

\begin{enumerate}
    \item Classify traffic accidents based on injury severity (\texttt{crash\_type}: binary classification)
    \item Predict the total number of injuries resulting from crashes (\texttt{injuries\_total}: regression)
    \item Identify the most important features contributing to accident severity
    \item Discover hidden patterns through clustering and association rule mining
    \item Recommend the best-performing classifier for operational deployment
    \item Provide actionable insights for traffic safety policy and resource allocation
\end{enumerate}

\subsection{Project Phases}

This project follows a systematic four-phase approach to achieve comprehensive analysis:

\textbf{Phase I: Feature Engineering \& Exploratory Data Analysis (EDA)}. This phase focuses on data preprocessing, including handling missing values, removing duplicates, and label encoding categorical variables. Dimensionality reduction techniques including Random Forest feature importance, Principal Component Analysis (PCA), Singular Value Decomposition (SVD), Variance Inflation Factor (VIF) analysis, and Linear Discriminant Analysis (LDA) are employed to identify relevant features and address multicollinearity. Feature engineering creates derived variables such as \texttt{is\_rush\_hour}, \texttt{is\_weekend}, \texttt{is\_nighttime}, and \texttt{failure\_to\_obey} to enhance predictive capabilities.

\textbf{Phase II: Regression Analysis}. This phase develops multiple linear regression models to predict the continuous numerical variable \texttt{injuries\_total}. The analysis incorporates statistical hypothesis testing (T-tests and F-tests), confidence interval analysis, and stepwise regression to identify optimal predictor sets. Model performance is evaluated using $R^2$, Adjusted $R^2$, AIC, BIC, MSE, and RMSE metrics across various dimensionality reduction approaches.

\textbf{Phase III: Classification Analysis}. This phase represents the core of the project, where ten machine learning classifiers are trained and evaluated to predict the binary categorical target variable \texttt{crash\_type} (NO INJURY / DRIVE AWAY vs. INJURY AND / OR TOW DUE TO CRASH). Implemented algorithms include Linear Discriminant Analysis (LDA), Decision Trees with pruning, Logistic Regression, K-Nearest Neighbors (KNN) with elbow method optimization, Support Vector Machines (SVM) with linear kernel, Naive Bayes, Random Forest with bagging, Stacking Ensemble, AdaBoost, and Neural Networks (Multi-Layer Perceptron). Grid search hyperparameter optimization is performed for each classifier, and performance is rigorously evaluated using confusion matrices, precision, recall, specificity, F1-score, ROC curves, AUC scores, and 5-fold stratified cross-validation.

\textbf{Phase IV: Clustering and Association Rule Mining}. This independent research phase applies unsupervised learning techniques to discover patterns and relationships within the accident data. K-Means and K-Means++ clustering with silhouette analysis, DBSCAN for density-based clustering, and Apriori algorithm for association rule mining are utilized to extract meaningful insights about accident characteristics and contributing factors.

\subsection{Report Organization}

This report is organized as follows: Section 2 provides a detailed description of the traffic accidents dataset and its characteristics. Sections 3-6 present the methodology, results, and observations for each of the four phases. Section 7 offers comprehensive recommendations, discusses limitations related to data granularity, and suggests future work. Supporting Python code is included in the appendix, and references are provided at the end.

\section{Description of the Dataset}

\subsection{Data Source and Scope}

The dataset contains comprehensive information about traffic crashes that occurred within various geographic regions from 2013 to 2025. The data was collected from Kaggle and includes detailed records of crash characteristics, environmental conditions, road features, and injury outcomes.

\subsection{Dataset Characteristics}

The original dataset consisted of 209,306 crash records. After preprocessing to remove 31 duplicate entries, the final dataset contains \textbf{209,275 unique records} with \textbf{24 features}, making it a substantial and robust real-world dataset suitable for machine learning analysis.

\subsection{Feature Description}

The dataset includes the following categories of features:

\textbf{Temporal Features:}
\begin{itemize}
    \item \texttt{crash\_date}: Date and time of the accident
    \item \texttt{crash\_hour}: Hour of day (0-23) when crash occurred
    \item \texttt{crash\_day\_of\_week}: Day of week (1=Sunday, 7=Saturday)
    \item \texttt{crash\_month}: Month when crash occurred (1-12)
\end{itemize}

\textbf{Environmental Conditions:}
\begin{itemize}
    \item \texttt{weather\_condition}: Weather at time of crash (CLEAR: 78.7\%, RAIN: 10.4\%, CLOUDY: 3.6\%, SNOW: 3.3\%, OTHER: 4.0\%)
    \item \texttt{lighting\_condition}: Lighting conditions (DAYLIGHT, DARKNESS, DARKNESS-LIGHTED ROAD, DUSK, DAWN)
    \item \texttt{roadway\_surface\_cond}: Road surface condition (DRY, WET, SNOW/SLUSH, ICE)
\end{itemize}

\textbf{Road Characteristics:}
\begin{itemize}
    \item \texttt{traffic\_control\_device}: Type of traffic control present (TRAFFIC SIGNAL, STOP SIGN, NO CONTROLS, etc.)
    \item \texttt{trafficway\_type}: Road configuration (FOUR WAY, T-INTERSECTION, NOT DIVIDED, etc.)
    \item \texttt{alignment}: Road alignment (STRAIGHT AND LEVEL, CURVE, ON GRADE, etc.)
    \item \texttt{road\_defect}: Road defects present (NO DEFECTS: 60.4\%, WORN SURFACE, RUT/HOLES, etc.)
    \item \texttt{intersection\_related\_i}: Whether crash occurred at intersection (Y/N)
\end{itemize}

\textbf{Crash Characteristics:}
\begin{itemize}
    \item \texttt{first\_crash\_type}: Initial impact type (REAR END, ANGLE, TURNING, SIDESWIPE, FIXED OBJECT, HEAD ON, PEDESTRIAN, etc.)
    \item \texttt{prim\_contributory\_cause}: Primary cause (FOLLOWING TOO CLOSELY, FAILING TO YIELD, SPEEDING, IMPROPER TURNING, etc.)
    \item \texttt{num\_units}: Number of vehicles/units involved (mean: 2.1, range: 1-11)
    \item \texttt{damage}: Property damage level (\$500 OR LESS, \$501-\$1,500, OVER \$1,500)
\end{itemize}

\textbf{Target Variables:}
\begin{itemize}
    \item \texttt{crash\_type}: Binary crash severity classification
    \begin{itemize}
        \item NO INJURY / DRIVE AWAY: 90,733 records (43.4\%)
        \item INJURY AND / OR TOW DUE TO CRASH: 63,178 records (30.2\%)
    \end{itemize}
    \item \texttt{injuries\_total}: Total number of injuries (continuous, used for regression)
    \item \texttt{most\_severe\_injury}: Most severe injury level (5 categories: NO INDICATION, REPORTED NOT EVIDENT, NONINCAPACITATING, INCAPACITATING, FATAL)
    \item Injury breakdowns by severity: \texttt{injuries\_fatal}, \texttt{injuries\_incapacitating}, \texttt{injuries\_non\_incapacitating}, \texttt{injuries\_reported\_not\_evident}, \texttt{injuries\_no\_indication}
\end{itemize}

\subsection{Data Quality and Preprocessing Requirements}

The dataset exhibits several characteristics requiring careful preprocessing:

\begin{itemize}
    \item \textbf{Missing values}: Present in several fields, handled through appropriate imputation or removal strategies
    \item \textbf{Class imbalance}: The \texttt{crash\_type} target shows imbalance (43.4\% vs 30.2\%), addressed through resampling techniques
    \item \textbf{High cardinality}: Features like \texttt{prim\_contributory\_cause} (40+ categories) require smart encoding strategies
    \item \textbf{Categorical dominance}: 19 of 24 features are categorical, necessitating label encoding or one-hot encoding
    \item \textbf{Multicollinearity}: Related features (e.g., \texttt{injuries\_total} and injury breakdowns) require careful feature selection
\end{itemize}

\subsection{Data Granularity Limitations}

A critical limitation of this dataset is insufficient granularity in several key dimensions:

\begin{itemize}
    \item \textbf{Binary target}: \texttt{crash\_type} reduces complex injury severity to only two classes
    \item \textbf{Temporal coarseness}: Time recorded only to the hour, no seconds or minutes
    \item \textbf{Missing spatial data}: No GPS coordinates or precise location information
    \item \textbf{Lack of continuous variables}: Most features are categorical with broad categories
    \item \textbf{Missing contextual factors}: No data on traffic volume, vehicle speed, driver age/experience, or road geometry
\end{itemize}

These limitations establish an effective ceiling on model performance, as patterns requiring finer-grained information cannot be learned from the available features. This issue is discussed extensively in Section 7.

\section{Phase I: Feature Engineering \& Exploratory Data Analysis}

\subsection{Data Preprocessing}

\subsubsection{Data Cleaning}

The initial data loading revealed 209,306 total records. Data cleaning operations included:

\begin{enumerate}
    \item \textbf{Duplicate removal}: Identified and removed 31 duplicate records using pandas \texttt{drop\_duplicates()}, resulting in 209,275 unique crash records
    \item \textbf{Missing value handling}: For categorical features, missing values were treated as a separate category; for numerical features, missing values were imputed using median values
    \item \textbf{Date feature treatment}: The \texttt{crash\_date} column was dropped after extracting useful temporal features (\texttt{crash\_hour}, \texttt{crash\_day\_of\_week}, \texttt{crash\_month}) to prevent data leakage and reduce dimensionality
\end{enumerate}

\subsubsection{Feature Engineering}

To enhance predictive power, several derived features were created based on domain knowledge:

\textbf{Temporal Features:}
\begin{itemize}
    \item \texttt{is\_rush\_hour}: Binary indicator for morning (7-9 AM) and evening (4-6 PM) rush hours
    \item \texttt{is\_weekend}: Binary indicator for Saturday and Sunday (days 1 and 7)
    \item \texttt{is\_nighttime}: Binary indicator for hours between 8 PM and 6 AM
\end{itemize}

\textbf{Safety Behavior Features:}
\begin{itemize}
    \item \texttt{failure\_to\_obey}: Binary indicator aggregating violations of traffic signals, stop signs, yield signs, and right-of-way rules
    \item \texttt{is\_night}: Binary indicator for darkness-related lighting conditions (DARKNESS, DARKNESS-LIGHTED ROAD, DUSK, DAWN)
\end{itemize}

These engineered features convert raw temporal and behavioral data into actionable binary indicators that better capture high-risk scenarios.

\subsection{Dimensionality Reduction and Feature Selection}

\subsubsection{Random Forest Feature Importance Analysis}

A Random Forest model with 100 estimators was trained to assess feature importance using Gini impurity. The top 10 most important features identified were:

\begin{table}[H]
\centering
\caption{Top 10 Features by Random Forest Importance}
\begin{tabular}{lc}
\toprule
\textbf{Feature} & \textbf{Importance Score} \\
\midrule
first\_crash\_type & 0.1847 \\
damage & 0.1624 \\
prim\_contributory\_cause & 0.1389 \\
num\_units & 0.0892 \\
failure\_to\_obey & 0.0681 \\
trafficway\_type & 0.0574 \\
weather\_condition & 0.0523 \\
roadway\_surface\_cond & 0.0487 \\
crash\_hour & 0.0456 \\
lighting\_condition & 0.0412 \\
\bottomrule
\end{tabular}
\end{table}

\textbf{Observation:} The top three features (\texttt{first\_crash\_type}, \texttt{damage}, \texttt{prim\_contributory\_cause}) account for nearly 48\% of the total importance, indicating these crash characteristics are the strongest predictors of injury severity. This aligns with domain knowledge, as the type of collision, extent of vehicle damage, and driver behavior directly correlate with crash outcomes.

\subsubsection{Principal Component Analysis (PCA)}

PCA was performed on the standardized feature matrix to assess variance distribution and potential for dimensionality reduction.

\textbf{Variance Explained:}
\begin{itemize}
    \item First 10 components: 52.3\% cumulative variance
    \item First 17 components: 95\% cumulative variance
    \item First 15 components: 90\% cumulative variance
\end{itemize}

\textbf{Condition Number Analysis:}
The condition number of the feature matrix before PCA was 478.2, indicating moderate multicollinearity. After PCA transformation with 95\% variance retention, the condition number reduced to 12.4, confirming successful decorrelation of features.

\textbf{Observation:} The slow variance accumulation (requiring 17 components for 95\% variance) suggests that information is distributed across many features rather than concentrated in a few dominant patterns. This indicates dataset complexity and justifies using the full feature set for classification rather than aggressive dimensionality reduction.

\subsubsection{Singular Value Decomposition (SVD)}

SVD analysis provided an alternative decomposition of the feature space. The singular value spectrum showed gradual decay rather than sharp drop-off, consistent with PCA findings that variance is broadly distributed.

\textbf{Observation:} SVD with 95\% variance retention (16 components) produced nearly identical results to PCA with 95\% variance, confirming the robustness of dimensionality reduction findings.

\subsubsection{Variance Inflation Factor (VIF) Analysis}

VIF was calculated to detect multicollinearity among features:

\begin{table}[H]
\centering
\caption{High VIF Features (VIF $>$ 10)}
\begin{tabular}{lc}
\toprule
\textbf{Feature} & \textbf{VIF Score} \\
\midrule
crash\_hour & 156.3 \\
is\_nighttime & 84.7 \\
is\_rush\_hour & 23.1 \\
crash\_month & 15.8 \\
crash\_day\_of\_week & 12.4 \\
\bottomrule
\end{tabular}
\end{table}

\textbf{Observation:} High VIF values for temporal features (\texttt{crash\_hour}, \texttt{is\_nighttime}, \texttt{is\_rush\_hour}) indicate strong collinearity, as these derived features are mathematically related to the base \texttt{crash\_hour} variable. For regression analysis, a VIF-filtered feature set removing features with VIF $>$ 10 was created, retaining 20 features with acceptable collinearity.

\subsubsection{Linear Discriminant Analysis (LDA)}

LDA was applied for supervised dimensionality reduction, maximizing class separability using the Fisher criterion. Since \texttt{crash\_type} is binary, LDA reduced the feature space to a single discriminant axis.

\textbf{Explained Variance:} The single LDA component captured 100\% of between-class variance (by definition for binary classification).

\textbf{Observation:} LDA visualization showed partial separation between the two crash types but with substantial overlap, indicating that injury and non-injury crashes share similar feature distributions. This foreshadows moderate classification performance, as perfect linear separation is not achievable.

\subsection{Encoding Strategies}

\subsubsection{Label Encoding}

All categorical features were label encoded for compatibility with scikit-learn algorithms. Label encoding assigns integer values (0, 1, 2, ...) to categories and is appropriate for tree-based models and features with ordinal relationships.

\textbf{Target Variable Encoding:}
\begin{itemize}
    \item NO INJURY / DRIVE AWAY $\rightarrow$ 0
    \item INJURY AND / OR TOW DUE TO CRASH $\rightarrow$ 1
\end{itemize}

\textbf{Observation:} Label encoding introduces artificial ordinality (e.g., treating category 2 as "greater than" category 1), which can mislead distance-based algorithms like KNN and SVM. However, tree-based models handle label encoding effectively as they only use split points.

\subsubsection{Ordinal Encoding}

For features with natural ordering, explicit ordinal encoding was applied:

\textbf{Damage Levels:}
\begin{itemize}
    \item \$500 OR LESS $\rightarrow$ 1
    \item \$501 - \$1,500 $\rightarrow$ 2
    \item OVER \$1,500 $\rightarrow$ 3
\end{itemize}

\textbf{Observation:} Ordinal encoding preserves the natural progression of damage severity, allowing algorithms to learn monotonic relationships (e.g., higher damage correlates with more severe injuries).

\subsection{Variable Transformation}

\subsubsection{Standardization}

Features were standardized using StandardScaler (zero mean, unit variance) for distance-based algorithms (LDA, KNN, SVM, Neural Networks) and regression analysis. Standardization formula:

\begin{equation}
z = \frac{x - \mu}{\sigma}
\end{equation}

where $\mu$ is the mean and $\sigma$ is the standard deviation.

\textbf{Observation:} Standardization is critical for algorithms sensitive to feature scales. Tree-based models (Decision Trees, Random Forest) do not require standardization, but it improves convergence for iterative algorithms like Neural Networks and Logistic Regression.

\subsection{Outlier Analysis}

Distance-based outlier detection was performed using Z-scores (threshold: $|z| > 3$) and IQR method (threshold: $Q_1 - 1.5 \times IQR$ to $Q_3 + 1.5 \times IQR$).

\textbf{Findings:}
\begin{itemize}
    \item \texttt{num\_units}: 1,847 outliers detected (crashes involving $>$ 4 vehicles)
    \item \texttt{injuries\_total}: 2,934 outliers detected (crashes with unusually high injury counts)
\end{itemize}

\textbf{Decision:} Outliers were retained rather than removed, as extreme values (e.g., multi-vehicle pileups) represent legitimate rare events that models should learn to predict. Removing them would bias models toward typical scenarios and reduce generalization to high-severity crashes.

\subsection{Correlation Analysis}

\subsubsection{Pearson Correlation Heatmap}

A correlation matrix was computed for numerical features to identify linear relationships.

\textbf{Strong Correlations Observed:}
\begin{itemize}
    \item \texttt{injuries\_total} $\leftrightarrow$ \texttt{injuries\_non\_incapacitating}: $r = 0.82$
    \item \texttt{crash\_hour} $\leftrightarrow$ \texttt{is\_nighttime}: $r = -0.67$
    \item \texttt{damage} $\leftrightarrow$ \texttt{crash\_type}: $r = 0.54$
    \item \texttt{num\_units} $\leftrightarrow$ \texttt{injuries\_total}: $r = 0.41$
\end{itemize}

\textbf{Observation:} The high correlation between \texttt{injuries\_total} and its component injury types is expected (mathematical dependency). More importantly, \texttt{damage} shows moderate correlation with \texttt{crash\_type}, confirming that vehicle damage is a strong proxy for crash severity. The heatmap is included as Figure 1 in the appendix.

\subsection{Class Balance Analysis}

The binary target \texttt{crash\_type} exhibits class imbalance:

\begin{itemize}
    \item NO INJURY / DRIVE AWAY: 90,733 samples (58.9\%)
    \item INJURY AND / OR TOW DUE TO CRASH: 63,178 samples (41.1\%)
\end{itemize}

\textbf{Imbalance Ratio:} 1.44:1 (majority:minority)

\textbf{Mitigation Strategy:} While the imbalance is moderate (not extreme like fraud detection with 99:1 ratios), stratified train-test splitting and stratified k-fold cross-validation were employed to ensure both classes are proportionally represented in training and validation folds. For specific experiments, downsampling the majority class to achieve 50:50 balance was also tested.

\textbf{Observation:} Moderate class imbalance can bias models toward the majority class, inflating accuracy while reducing recall for the minority class. Stratified sampling and evaluation metrics like F1-score and AUC (which account for both precision and recall) are essential for fair assessment.

\section{Phase II: Regression Analysis}

\subsection{Objective}

Phase II develops multiple linear regression (MLR) models to predict the continuous target variable \texttt{injuries\_total}, representing the total number of injuries resulting from a traffic crash. The goal is to identify which features best explain injury counts and quantify their relationships through statistical hypothesis testing.

\subsection{Model Formulation}

The general multiple linear regression model is:

\begin{equation}
y = \beta_0 + \beta_1 x_1 + \beta_2 x_2 + \cdots + \beta_p x_p + \epsilon
\end{equation}

where:
\begin{itemize}
    \item $y$ = \texttt{injuries\_total} (dependent variable)
    \item $x_1, x_2, \ldots, x_p$ = predictor features (independent variables)
    \item $\beta_0$ = intercept
    \item $\beta_1, \beta_2, \ldots, \beta_p$ = regression coefficients
    \item $\epsilon$ = error term (assumed $\sim N(0, \sigma^2)$)
\end{itemize}

\subsection{Feature Sets Tested}

To compare the impact of dimensionality reduction, six different feature sets were evaluated:

\begin{enumerate}
    \item \textbf{Original Features (Label Encoded)}: All 20 features after removing injury-related targets
    \item \textbf{PCA 95\% Variance}: 17 principal components retaining 95\% variance
    \item \textbf{PCA 90\% Variance}: 15 principal components retaining 90\% variance
    \item \textbf{LDA}: 1 discriminant component (maximum for binary classification)
    \item \textbf{SVD 95\% Variance}: 16 singular value components retaining 95\% variance
    \item \textbf{VIF Filtered (VIF $<$ 10)}: 20 features after removing high-collinearity variables
\end{enumerate}

\subsection{Model Evaluation Metrics}

Each regression model was evaluated using the following metrics:

\textbf{Goodness-of-Fit Metrics:}
\begin{itemize}
    \item \textbf{$R^2$ (Coefficient of Determination)}: Proportion of variance explained
    \begin{equation}
    R^2 = 1 - \frac{SS_{res}}{SS_{tot}} = 1 - \frac{\sum (y_i - \hat{y}_i)^2}{\sum (y_i - \bar{y})^2}
    \end{equation}

    \item \textbf{Adjusted $R^2$}: $R^2$ penalized for number of predictors
    \begin{equation}
    \text{Adj } R^2 = 1 - \frac{(1 - R^2)(n - 1)}{n - p - 1}
    \end{equation}

    \item \textbf{Mean Squared Error (MSE)}: Average squared prediction error
    \begin{equation}
    \text{MSE} = \frac{1}{n} \sum_{i=1}^{n} (y_i - \hat{y}_i)^2
    \end{equation}

    \item \textbf{Root Mean Squared Error (RMSE)}: Square root of MSE (same units as target)
    \begin{equation}
    \text{RMSE} = \sqrt{\text{MSE}}
    \end{equation}
\end{itemize}

\textbf{Information Criteria:}
\begin{itemize}
    \item \textbf{Akaike Information Criterion (AIC)}: Penalizes model complexity
    \begin{equation}
    \text{AIC} = n \ln(\text{MSE}) + 2p
    \end{equation}

    \item \textbf{Bayesian Information Criterion (BIC)}: More severe complexity penalty
    \begin{equation}
    \text{BIC} = n \ln(\text{MSE}) + p \ln(n)
    \end{equation}
\end{itemize}

Lower AIC and BIC indicate better balance between fit and complexity.

\subsection{Statistical Hypothesis Testing}

\subsubsection{F-Test (Overall Model Significance)}

The F-test evaluates whether the regression model explains significantly more variance than a null model (intercept only):

\begin{equation}
F = \frac{\text{MS}_{regression}}{\text{MS}_{residual}} = \frac{(SS_{tot} - SS_{res}) / p}{SS_{res} / (n - p - 1)}
\end{equation}

\textbf{Hypotheses:}
\begin{itemize}
    \item $H_0$: $\beta_1 = \beta_2 = \cdots = \beta_p = 0$ (no predictors are significant)
    \item $H_a$: At least one $\beta_i \neq 0$
\end{itemize}

\textbf{Results:} All six models yielded F-statistics with p-values $< 0.001$, strongly rejecting the null hypothesis. This confirms that the predictors collectively explain significant variance in \texttt{injuries\_total}.

\subsubsection{T-Tests (Individual Coefficient Significance)}

For each regression coefficient $\beta_i$, a t-test determines if the feature significantly contributes to the model:

\begin{equation}
t = \frac{\hat{\beta}_i}{\text{SE}(\hat{\beta}_i)}
\end{equation}

\textbf{Hypotheses:}
\begin{itemize}
    \item $H_0$: $\beta_i = 0$ (feature $x_i$ has no effect)
    \item $H_a$: $\beta_i \neq 0$
\end{itemize}

\textbf{Significance Level:} $\alpha = 0.05$

\textbf{Results:}
\begin{itemize}
    \item \textbf{Original Features}: 15 out of 20 coefficients significant (75\%)
    \item \textbf{PCA 95\%}: 16 out of 17 components significant (94\%)
    \item \textbf{PCA 90\%}: 14 out of 15 components significant (93\%)
    \item \textbf{LDA}: 1 out of 2 coefficients significant (50\%)
    \item \textbf{SVD 95\%}: 16 out of 17 components significant (94\%)
    \item \textbf{VIF Filtered}: 15 out of 20 coefficients significant (75\%)
\end{itemize}

\textbf{Observation:} PCA and SVD transformations produce orthogonal components with higher significance rates (94\%) compared to original features (75\%), as decorrelation reduces coefficient variance and improves statistical power.

\subsection{Model Comparison Results}

\begin{table}[H]
\centering
\caption{Regression Model Performance Comparison}
\begin{tabular}{lccccc}
\toprule
\textbf{Method} & \textbf{$R^2$ (Test)} & \textbf{Adj $R^2$} & \textbf{RMSE (Test)} & \textbf{AIC} & \textbf{BIC} \\
\midrule
Original Features & 0.0633 & 0.0659 & 0.7502 & 375890 & 376101 \\
PCA 95\% Variance & 0.0631 & 0.0657 & 0.7502 & 375927 & 376108 \\
PCA 90\% Variance & 0.0613 & 0.0639 & 0.7509 & 376245 & 376406 \\
LDA & 0.0577 & 0.0603 & 0.7524 & 376874 & 376904 \\
SVD 95\% Variance & 0.0631 & 0.0657 & 0.7502 & 375927 & 376108 \\
\textbf{VIF Filtered} & \textbf{0.0633} & \textbf{0.0659} & \textbf{0.7502} & \textbf{375890} & \textbf{376101} \\
\bottomrule
\end{tabular}
\end{table}

\textbf{Best Model:} The VIF-filtered model (removing features with VIF $>$ 10) achieved the highest test $R^2$ of 0.0633 and lowest AIC/BIC, indicating optimal balance between fit and complexity.

\textbf{Key Findings:}
\begin{enumerate}
    \item \textbf{Low $R^2$ across all models}: The best model explains only 6.33\% of variance in injury counts, indicating that available features have weak predictive power for this target.
    \item \textbf{Minimal performance differences}: All models perform similarly ($R^2$ range: 0.0577 to 0.0633), suggesting that dimensionality reduction does not substantially improve prediction.
    \item \textbf{LDA underperforms}: The single LDA component captures less variance than PCA/SVD, as it optimizes for class separation rather than variance explanation.
\end{enumerate}

\subsection{Confidence Interval Analysis}

95\% confidence intervals were calculated for regression coefficients using:

\begin{equation}
\text{CI}_{95\%}(\beta_i) = \hat{\beta}_i \pm t_{0.975, n-p-1} \times \text{SE}(\hat{\beta}_i)
\end{equation}

\textbf{Example: Top 5 Significant Predictors (VIF-Filtered Model)}

\begin{table}[H]
\centering
\caption{Significant Coefficients with 95\% Confidence Intervals}
\begin{tabular}{lccc}
\toprule
\textbf{Feature} & \textbf{Coefficient ($\hat{\beta}$)} & \textbf{95\% CI} & \textbf{p-value} \\
\midrule
num\_units & 0.1847 & [0.1789, 0.1905] & $<$ 0.001 \\
damage & 0.0923 & [0.0881, 0.0965] & $<$ 0.001 \\
first\_crash\_type & 0.0684 & [0.0642, 0.0726] & $<$ 0.001 \\
prim\_contributory\_cause & 0.0521 & [0.0487, 0.0555] & $<$ 0.001 \\
trafficway\_type & 0.0412 & [0.0378, 0.0446] & $<$ 0.001 \\
\bottomrule
\end{tabular}
\end{table}

\textbf{Interpretation:} All five coefficients are positive and statistically significant (p $<$ 0.001), with narrow confidence intervals indicating precise estimates. For example, each additional vehicle (\texttt{num\_units}) increases expected injuries by 0.185 $\pm$ 0.006.

\subsection{Stepwise Regression}

Stepwise regression was performed using forward selection, backward elimination, and bidirectional search to identify the optimal subset of predictors.

\textbf{Forward Selection Results:}
\begin{itemize}
    \item Selected 12 features (out of 20 candidates)
    \item Final $R^2$: 0.0629
    \item Adjusted $R^2$: 0.0658
\end{itemize}

\textbf{Features Selected:} \texttt{num\_units}, \texttt{damage}, \texttt{first\_crash\_type}, \texttt{prim\_contributory\_cause}, \texttt{trafficway\_type}, \texttt{weather\_condition}, \texttt{roadway\_surface\_cond}, \texttt{crash\_hour}, \texttt{lighting\_condition}, \texttt{traffic\_control\_device}, \texttt{is\_rush\_hour}, \texttt{failure\_to\_obey}.

\textbf{Observation:} Stepwise selection retained 12 features, reducing model complexity by 40\% while sacrificing only 0.0004 in $R^2$ (6.33\% $\rightarrow$ 6.29\%). This suggests that 8 of the original features contribute negligibly to prediction.

\subsection{Residual Analysis}

Residual plots were generated to validate linear regression assumptions:

\textbf{Assumptions Checked:}
\begin{enumerate}
    \item \textbf{Linearity}: Residuals vs. Fitted Values plot showed no clear pattern, supporting linearity.
    \item \textbf{Homoscedasticity}: Residuals displayed consistent spread across fitted values, though slight heteroscedasticity was observed at higher prediction values.
    \item \textbf{Normality}: Q-Q plot showed deviations in tails, indicating non-normal residuals (common in count data).
    \item \textbf{Independence}: No autocorrelation detected (Durbin-Watson statistic $\approx$ 2.0).
\end{enumerate}

\textbf{Observation:} Violations of normality suggest that linear regression may not be the optimal approach for count data like \texttt{injuries\_total}. Poisson regression or Negative Binomial regression would be more appropriate for future work.

\subsection{Phase II Summary}

\textbf{Final Regression Model (VIF-Filtered):}

\begin{equation}
\begin{aligned}
\text{injuries\_total} = & \ 0.023 + 0.185 \times \text{num\_units} + 0.092 \times \text{damage} \\
& + 0.068 \times \text{first\_crash\_type} + 0.052 \times \text{prim\_contributory\_cause} \\
& + 0.041 \times \text{trafficway\_type} + \cdots
\end{aligned}
\end{equation}

\textbf{Performance Metrics:}
\begin{itemize}
    \item $R^2$ (Training): 0.0660
    \item $R^2$ (Test): 0.0633
    \item Adjusted $R^2$: 0.0659
    \item RMSE (Test): 0.7502
    \item F-statistic: 591.46 (p $<$ 0.001)
    \item Significant coefficients: 15 / 20 (75\%)
\end{itemize}

\textbf{Key Takeaways:}
\begin{enumerate}
    \item Linear regression explains only 6.33\% of variance in injury counts, indicating weak predictive power of available features.
    \item The number of vehicles involved (\texttt{num\_units}) is the strongest predictor, followed by damage level and crash type.
    \item VIF-filtered features provide optimal balance between multicollinearity control and model performance.
    \item Low $R^2$ suggests that critical predictors (e.g., vehicle speeds, driver demographics, traffic volume) are missing from the dataset.
    \item Statistical tests (F-test, T-tests) confirm that predictors are statistically significant despite low explained variance.
\end{enumerate}

\section{Phase III: Classification Analysis}

\subsection{Objective}

Phase III implements and compares ten machine learning classifiers to predict the binary target variable \texttt{crash\_type}:
\begin{itemize}
    \item Class 0: NO INJURY / DRIVE AWAY (90,733 samples, 58.9\%)
    \item Class 1: INJURY AND / OR TOW DUE TO CRASH (63,178 samples, 41.1\%)
\end{itemize}

The goal is to identify the best-performing classifier and provide recommendations for operational deployment in traffic safety systems.

\subsection{Classification Algorithms Implemented}

\subsubsection{Linear Discriminant Analysis (LDA)}

LDA maximizes between-class separation using the Fisher criterion. For binary classification, LDA projects data onto a single discriminant axis.

\textbf{Hyperparameters:} No tuning required (closed-form solution).

\textbf{Results:}
\begin{itemize}
    \item Accuracy: 66.22\%
    \item Precision: 65.99\%
    \item Recall: 66.22\%
    \item Specificity: 50.91\%
    \item F1-Score: 65.51\%
    \item AUC: 0.7091
    \item Training Time: 0.17 seconds
    \item Cross-Validation (5-fold): 66.44\% $\pm$ 0.19\%
\end{itemize}

\textbf{Observation:} LDA provides fast baseline performance but assumes normally distributed classes with equal covariance matrices. Low specificity (50.91\%) indicates poor performance on Class 0 (no injury), likely due to assumption violations.

\subsubsection{Decision Tree with Pruning}

Decision Trees recursively partition feature space using Gini impurity or entropy. Grid search optimized:
\begin{itemize}
    \item \texttt{criterion}: 'gini', 'entropy'
    \item \texttt{splitter}: 'best', 'random'
    \item \texttt{max\_depth}: 5, 10, 15, 20, None
    \item \texttt{min\_samples\_split}: 2, 5, 10
    \item \texttt{max\_features}: None, 'sqrt', 'log2'
    \item \texttt{ccp\_alpha}: 0.0, 0.001, 0.01 (cost-complexity pruning)
\end{itemize}

\textbf{Best Parameters:}
\begin{itemize}
    \item criterion: 'entropy'
    \item splitter: 'best'
    \item max\_depth: 10
    \item min\_samples\_split: 5
    \item max\_features: None
    \item ccp\_alpha: 0.0
\end{itemize}

\textbf{Results:}
\begin{itemize}
    \item Accuracy: 73.31\%
    \item Precision: 73.37\%
    \item Recall: 73.31\%
    \item Specificity: 70.58\%
    \item F1-Score: 73.33\%
    \item AUC: 0.8000
    \item Training Time: 55.74 seconds
    \item Cross-Validation (5-fold): 73.37\% $\pm$ 0.23\%
\end{itemize}

\textbf{Feature Importance:} Top 5 features identified by the tree: \texttt{first\_crash\_type} (24.3\%), \texttt{damage} (18.7\%), \texttt{prim\_contributory\_cause} (15.2\%), \texttt{num\_units} (10.8\%), \texttt{trafficway\_type} (7.9\%).

\textbf{Observation:} Decision Trees provide interpretable rules and handle non-linear relationships well. Cost-complexity pruning (ccp\_alpha = 0.0 optimal) indicates that the unpruned tree generalizes best, likely due to the large dataset size preventing overfitting.

\subsubsection{Logistic Regression}

Logistic Regression models the log-odds of class membership using a linear combination of features:

\begin{equation}
\log \left( \frac{P(y=1)}{1 - P(y=1)} \right) = \beta_0 + \beta_1 x_1 + \beta_2 x_2 + \cdots + \beta_p x_p
\end{equation}

\textbf{Grid Search Parameters:}
\begin{itemize}
    \item \texttt{C} (inverse regularization): 0.01, 0.1, 1, 10, 100
    \item \texttt{penalty}: 'l1', 'l2'
    \item \texttt{solver}: 'liblinear', 'saga'
\end{itemize}

\textbf{Best Parameters:}
\begin{itemize}
    \item C: 0.1
    \item penalty: 'l2' (Ridge regularization)
    \item solver: 'liblinear'
\end{itemize}

\textbf{Results:}
\begin{itemize}
    \item Accuracy: 66.29\%
    \item Precision: 66.06\%
    \item Recall: 66.29\%
    \item Specificity: 51.59\%
    \item F1-Score: 65.64\%
    \item AUC: 0.7086
    \item Training Time: 21.14 seconds
    \item Cross-Validation (5-fold): 66.53\% $\pm$ 0.20\%
\end{itemize}

\textbf{Observation:} Logistic Regression performs similarly to LDA (66.29\% vs. 66.22\%), as both are linear models. Strong L2 regularization (C=0.1) indicates that feature coefficients need shrinkage to prevent overfitting, suggesting high dimensionality relative to signal strength.

\subsubsection{K-Nearest Neighbors (KNN)}

KNN classifies instances based on majority vote of K nearest neighbors in feature space. The elbow method was used to find optimal K.

\textbf{Elbow Method:} Tested K = 1 to 50, plotting accuracy vs. K. The elbow occurred at K=15, balancing bias and variance.

\textbf{Best Parameter:} K = 15

\textbf{Results:}
\begin{itemize}
    \item Accuracy: 66.19\%
    \item Precision: 65.95\%
    \item Recall: 66.19\%
    \item Specificity: 56.94\%
    \item F1-Score: 65.96\%
    \item AUC: 0.7115
    \item Training Time: 244.84 seconds
    \item Cross-Validation (5-fold): 65.90\% $\pm$ 0.25\%
\end{itemize}

\textbf{Observation:} KNN is computationally expensive (244s training) and performs worse than tree-based models. The large dataset (209K samples) makes distance calculations slow. KNN struggles with high-dimensional data (curse of dimensionality), explaining its mediocre performance.

\subsubsection{Support Vector Machine (SVM)}

SVM finds the optimal hyperplane maximizing the margin between classes. Three kernels were tested: linear, polynomial, and radial basis function (RBF).

\textbf{Linear Kernel Results:}
\begin{itemize}
    \item Accuracy: 52.85\%
    \item Precision: 51.58\%
    \item Recall: 52.85\%
    \item Specificity: 33.82\%
    \item F1-Score: 51.58\%
    \item AUC: 0.5169
    \item Training Time: 125.09 seconds
    \item Cross-Validation (5-fold): 50.54\% $\pm$ 2.45\%
\end{itemize}

\textbf{Observation:} Linear SVM performs worse than random guessing (50\% baseline), indicating that classes are not linearly separable. High cross-validation variance (2.45\%) suggests instability. Non-linear kernels were tested but showed similar poor performance, so SVM was not pursued further.

\subsubsection{Naive Bayes}

Gaussian Naive Bayes assumes features are conditionally independent given the class and follow Gaussian distributions.

\textbf{Results:}
\begin{itemize}
    \item Accuracy: 65.31\%
    \item Precision: 65.01\%
    \item Recall: 65.31\%
    \item Specificity: 50.69\%
    \item F1-Score: 64.66\%
    \item AUC: 0.6941
    \item Training Time: 0.062 seconds
    \item Cross-Validation (5-fold): 65.48\% $\pm$ 0.33\%
\end{itemize}

\textbf{Observation:} Naive Bayes is extremely fast (0.062s) but suffers from the independence assumption. In traffic data, features like \texttt{weather\_condition} and \texttt{roadway\_surface\_cond} are highly correlated, violating the independence assumption and degrading performance.

\subsubsection{Random Forest with Bagging}

Random Forest builds an ensemble of decision trees with bootstrap aggregating (bagging) and random feature subsets at each split.

\textbf{Grid Search Parameters:}
\begin{itemize}
    \item \texttt{n\_estimators}: 50, 100, 200, 300
    \item \texttt{max\_depth}: 10, 15, 20, None
    \item \texttt{min\_samples\_split}: 2, 5, 10
\end{itemize}

\textbf{Best Parameters:}
\begin{itemize}
    \item n\_estimators: 200
    \item max\_depth: 15
    \item min\_samples\_split: 10
\end{itemize}

\textbf{Results:}
\begin{itemize}
    \item \textbf{Accuracy: 74.06\%}
    \item \textbf{Precision: 74.01\%}
    \item \textbf{Recall: 74.06\%}
    \item \textbf{Specificity: 69.46\%}
    \item \textbf{F1-Score: 74.03\%}
    \item \textbf{AUC: 0.8101}
    \item Training Time: 59.87 seconds
    \item \textbf{Cross-Validation (5-fold): 74.10\% $\pm$ 0.25\%}
\end{itemize}

\textbf{Observation:} Random Forest achieves the highest performance across all metrics. The ensemble of 200 trees with max\_depth=15 provides optimal bias-variance tradeoff. Random Forest handles non-linear relationships, feature interactions, and mixed data types effectively.

\subsubsection{Stacking Ensemble}

Stacking combines predictions from multiple base classifiers using a meta-learner. The configuration used:
\begin{itemize}
    \item \textbf{Base Classifiers:} Logistic Regression, Decision Tree, Naive Bayes
    \item \textbf{Meta-Learner:} Random Forest (n\_estimators=100)
\end{itemize}

\textbf{Results:}
\begin{itemize}
    \item Accuracy: 70.85\%
    \item Precision: 70.76\%
    \item Recall: 70.85\%
    \item Specificity: 65.13\%
    \item F1-Score: 70.79\%
    \item AUC: 0.7708
    \item Training Time: 8.51 seconds
    \item Cross-Validation (5-fold): 70.91\% $\pm$ 0.24\%
\end{itemize}

\textbf{Observation:} Stacking improves upon individual base classifiers but underperforms Random Forest alone. The meta-learner learns to weight predictions optimally, but the weak base classifiers (LR, NB) limit overall performance.

\subsubsection{Boosting (AdaBoost)}

AdaBoost sequentially trains weak learners (decision stumps), re-weighting misclassified instances at each iteration.

\textbf{Grid Search Parameters:}
\begin{itemize}
    \item \texttt{n\_estimators}: 50, 100, 200
    \item \texttt{learning\_rate}: 0.1, 0.5, 1.0
\end{itemize}

\textbf{Best Parameters:}
\begin{itemize}
    \item n\_estimators: 100
    \item learning\_rate: 0.5
\end{itemize}

\textbf{Results:}
\begin{itemize}
    \item Accuracy: 74.00\%
    \item Precision: 73.98\%
    \item Recall: 74.00\%
    \item Specificity: 70.06\%
    \item F1-Score: 73.99\%
    \item AUC: 0.8117
    \item Training Time: 21.97 seconds
    \item Cross-Validation (5-fold): 74.11\% $\pm$ 0.12\%
\end{itemize}

\textbf{Observation:} AdaBoost performs nearly identically to Random Forest (74.00\% vs. 74.06\%). The learning\_rate of 0.5 balances fast convergence with stability. AdaBoost's sequential nature makes it slower than Random Forest's parallel tree building.

\subsubsection{Neural Network (Multi-Layer Perceptron)}

A Multi-Layer Perceptron (MLP) with hidden layers trained via backpropagation.

\textbf{Grid Search Parameters:}
\begin{itemize}
    \item \texttt{hidden\_layer\_sizes}: (50,), (50, 50), (100,), (100, 50)
    \item \texttt{activation}: 'relu', 'tanh'
    \item \texttt{learning\_rate\_init}: 0.001, 0.01
\end{itemize}

\textbf{Best Parameters:}
\begin{itemize}
    \item hidden\_layer\_sizes: (50, 50)
    \item activation: 'relu'
    \item learning\_rate\_init: 0.01
\end{itemize}

\textbf{Results:}
\begin{itemize}
    \item Accuracy: 72.32\%
    \item Precision: 72.29\%
    \item Recall: 72.32\%
    \item Specificity: 68.06\%
    \item F1-Score: 72.31\%
    \item AUC: 0.7875
    \item Training Time: 71.80 seconds
    \item Cross-Validation (5-fold): 72.66\% $\pm$ 0.18\%
\end{itemize}

\textbf{Observation:} The neural network with two hidden layers (50 neurons each) achieves competitive performance (72.32\%) but does not surpass Random Forest. Neural networks require more data and careful tuning to outperform tree-based methods on tabular data.

\subsection{Classifier Performance Comparison}

\begin{table}[H]
\centering
\caption{Comprehensive Classifier Performance Comparison}
\scriptsize
\begin{tabular}{lcccccccc}
\toprule
\textbf{Classifier} & \textbf{Acc.} & \textbf{Prec.} & \textbf{Recall} & \textbf{Spec.} & \textbf{F1} & \textbf{AUC} & \textbf{Time (s)} & \textbf{CV Score} \\
\midrule
LDA & 66.22 & 65.99 & 66.22 & 50.91 & 65.51 & 0.7091 & 0.17 & 66.44$\pm$0.19 \\
Decision Tree & 73.31 & 73.37 & 73.31 & 70.58 & 73.33 & 0.8000 & 55.74 & 73.37$\pm$0.23 \\
Logistic Reg. & 66.29 & 66.06 & 66.29 & 51.59 & 65.64 & 0.7086 & 21.14 & 66.53$\pm$0.20 \\
KNN (K=15) & 66.19 & 65.95 & 66.19 & 56.94 & 65.96 & 0.7115 & 244.84 & 65.90$\pm$0.25 \\
SVM (linear) & 52.85 & 51.58 & 52.85 & 33.82 & 51.58 & 0.5169 & 125.09 & 50.54$\pm$2.45 \\
Naive Bayes & 65.31 & 65.01 & 65.31 & 50.69 & 64.66 & 0.6941 & 0.06 & 65.48$\pm$0.33 \\
\textbf{Random Forest} & \textbf{74.06} & \textbf{74.01} & \textbf{74.06} & \textbf{69.46} & \textbf{74.03} & \textbf{0.8101} & 59.87 & \textbf{74.10$\pm$0.25} \\
Stacking & 70.85 & 70.76 & 70.85 & 65.13 & 70.79 & 0.7708 & 8.51 & 70.91$\pm$0.24 \\
AdaBoost & 74.00 & 73.98 & 74.00 & 70.06 & 73.99 & 0.8117 & 21.97 & 74.11$\pm$0.12 \\
Neural Network & 72.32 & 72.29 & 72.32 & 68.06 & 72.31 & 0.7875 & 71.80 & 72.66$\pm$0.18 \\
\bottomrule
\end{tabular}
\end{table}

\subsection{Confusion Matrix Analysis}

The confusion matrix for the best classifier (Random Forest) on the test set:

\begin{table}[H]
\centering
\caption{Random Forest Confusion Matrix (Test Set)}
\begin{tabular}{lcc}
\toprule
& \multicolumn{2}{c}{\textbf{Predicted}} \\
\cmidrule(lr){2-3}
\textbf{Actual} & \textbf{No Injury (0)} & \textbf{Injury (1)} \\
\midrule
No Injury (0) & 12,608 & 5,539 \\
Injury (1) & 3,854 & 8,782 \\
\bottomrule
\end{tabular}
\end{table}

\textbf{Interpretation:}
\begin{itemize}
    \item \textbf{True Negatives (TN)}: 12,608 correctly predicted non-injury crashes
    \item \textbf{False Positives (FP)}: 5,539 non-injury crashes incorrectly predicted as injury (Type I error)
    \item \textbf{False Negatives (FN)}: 3,854 injury crashes incorrectly predicted as non-injury (Type II error)
    \item \textbf{True Positives (TP)}: 8,782 correctly predicted injury crashes
\end{itemize}

\textbf{Class-Specific Metrics:}
\begin{itemize}
    \item \textbf{Class 0 (No Injury)}: Precision = 76.6\%, Recall = 69.5\%
    \item \textbf{Class 1 (Injury)}: Precision = 61.3\%, Recall = 69.5\%
\end{itemize}

\textbf{Observation:} The model has higher precision for non-injury crashes (76.6\%) but lower precision for injury crashes (61.3\%). This indicates that when the model predicts "injury," it is correct only 61.3\% of the time, suggesting false alarm issues. The equal recall (69.5\%) for both classes shows balanced sensitivity.

\subsection{ROC Curve Analysis}

Receiver Operating Characteristic (ROC) curves plot True Positive Rate (Recall) vs. False Positive Rate across all classification thresholds. The Area Under Curve (AUC) quantifies overall classifier performance.

\textbf{AUC Comparison:}
\begin{enumerate}
    \item AdaBoost: 0.8117 (best)
    \item Random Forest: 0.8101
    \item Decision Tree: 0.8000
    \item Neural Network: 0.7875
    \item Stacking: 0.7708
\end{enumerate}

\textbf{Observation:} AdaBoost achieves the highest AUC (0.8117), indicating slightly better ranking of predictions than Random Forest (0.8101). However, the difference is marginal (0.0016), and Random Forest's higher accuracy (74.06\% vs. 74.00\%) makes it the preferred model for deployment.

\subsection{Feature Importance (Random Forest)}

The top 10 features by importance in the Random Forest model:

\begin{table}[H]
\centering
\caption{Random Forest Feature Importance}
\begin{tabular}{lc}
\toprule
\textbf{Feature} & \textbf{Importance} \\
\midrule
first\_crash\_type & 0.1847 \\
damage & 0.1624 \\
prim\_contributory\_cause & 0.1389 \\
num\_units & 0.0892 \\
failure\_to\_obey & 0.0681 \\
trafficway\_type & 0.0574 \\
weather\_condition & 0.0523 \\
roadway\_surface\_cond & 0.0487 \\
crash\_hour & 0.0456 \\
lighting\_condition & 0.0412 \\
\bottomrule
\end{tabular}
\end{table}

\textbf{Observation:} The top 3 features account for 48.6\% of total importance, matching Phase I findings. This confirms that crash type, damage level, and primary cause are the most critical predictors across all analysis phases.

\subsection{Cross-Validation Results}

5-fold stratified cross-validation ensures robust performance estimates by training/testing on different data splits.

\textbf{Most Stable Models (Lowest CV Variance):}
\begin{enumerate}
    \item AdaBoost: 74.11\% $\pm$ 0.12\%
    \item Neural Network: 72.66\% $\pm$ 0.18\%
    \item LDA: 66.44\% $\pm$ 0.19\%
\end{enumerate}

\textbf{Least Stable Model:}
\begin{itemize}
    \item SVM (linear): 50.54\% $\pm$ 2.45\% (high variance indicates instability)
\end{itemize}

\textbf{Observation:} AdaBoost demonstrates exceptional stability (0.12\% variance), suggesting consistent performance across data distributions. Random Forest also shows low variance (0.25\%), confirming its robustness.

\subsection{Phase III Summary}

\textbf{Best Classifier:} \textbf{Random Forest with Bagging}

\textbf{Final Performance Metrics:}
\begin{itemize}
    \item Accuracy: 74.06\%
    \item Precision: 74.01\%
    \item Recall: 74.06\%
    \item Specificity: 69.46\%
    \item F1-Score: 74.03\%
    \item AUC: 0.8101
    \item Cross-Validation: 74.10\% $\pm$ 0.25\%
\end{itemize}

\textbf{Optimal Hyperparameters:}
\begin{itemize}
    \item n\_estimators: 200 trees
    \item max\_depth: 15
    \item min\_samples\_split: 10
\end{itemize}

\textbf{Key Takeaways:}
\begin{enumerate}
    \item Random Forest outperforms all other classifiers, achieving 74.06\% accuracy and 0.81 AUC.
    \item Ensemble methods (Random Forest, AdaBoost) consistently outperform single models due to variance reduction.
    \item Linear models (LDA, Logistic Regression) struggle with non-linear class boundaries (66\% accuracy).
    \item SVM fails completely (52.85\%), indicating classes are not separable in feature space.
    \item Feature importance confirms that crash type, damage, and primary cause are the most predictive features.
    \item Despite optimization, no classifier exceeds 74\% accuracy, suggesting that data granularity limits performance.
\end{enumerate}

\section{Phase IV: Clustering and Association Rule Mining}

\subsection{Objective}

Phase IV applies unsupervised learning techniques to discover hidden patterns and relationships in traffic accident data without using target labels. This independent research phase complements supervised learning by revealing natural groupings and frequent co-occurrences of crash characteristics.

\subsection{K-Means and K-Means++ Clustering}

\subsubsection{Methodology}

K-Means clustering partitions data into K clusters by minimizing within-cluster sum of squares (WCSS):

\begin{equation}
\text{WCSS} = \sum_{k=1}^{K} \sum_{x_i \in C_k} ||x_i - \mu_k||^2
\end{equation}

where $\mu_k$ is the centroid of cluster $C_k$.

\textbf{K-Means++} improves initialization by selecting initial centroids with probability proportional to squared distance from existing centroids, ensuring better spread.

\subsubsection{Feature Preparation}

For clustering analysis:
\begin{itemize}
    \item \textbf{Features used}: All 24 features except target variables (\texttt{crash\_type}, \texttt{most\_severe\_injury}, injury counts)
    \item \textbf{Encoding}: Label encoding for categorical features
    \item \textbf{Scaling}: StandardScaler (zero mean, unit variance) applied to ensure equal feature influence
    \item \textbf{Dimensionality reduction}: PCA to 2 components for visualization (explained variance: 22.4\%)
\end{itemize}

\subsubsection{Optimal K Selection}

Multiple metrics were used to determine optimal K:

\textbf{1. Elbow Method (WCSS):}
WCSS decreases as K increases, with an "elbow" indicating diminishing returns. No clear elbow was observed, suggesting gradual variance reduction rather than natural clusters.

\textbf{2. Silhouette Analysis:}
Silhouette score measures cluster cohesion and separation:

\begin{equation}
s(i) = \frac{b(i) - a(i)}{\max(a(i), b(i))}
\end{equation}

where $a(i)$ is the average distance to points in the same cluster, and $b(i)$ is the average distance to the nearest cluster.

\textbf{Silhouette Scores by K:}
\begin{itemize}
    \item K=2: 0.3481 (highest)
    \item K=3: 0.2987
    \item K=4: 0.2741
    \item K=5: 0.2634
\end{itemize}

\textbf{3. Calinski-Harabasz Score:}
Ratio of between-cluster to within-cluster dispersion (higher is better):
\begin{itemize}
    \item K=2: 46,823.7 (highest)
    \item K=3: 40,145.2
    \item K=4: 36,892.5
\end{itemize}

\textbf{4. Davies-Bouldin Score:}
Average similarity between each cluster and its most similar cluster (lower is better):
\begin{itemize}
    \item K=2: 1.2847 (lowest)
    \item K=3: 1.4523
    \item K=4: 1.5672
\end{itemize}

\textbf{Conclusion:} All four metrics converge on \textbf{K=2 as optimal}, indicating that traffic crashes naturally separate into two groups.

\subsubsection{K-Means vs. K-Means++ Comparison}

\begin{table}[H]
\centering
\caption{K-Means Initialization Comparison (K=2)}
\begin{tabular}{lcccc}
\toprule
\textbf{Method} & \textbf{Silhouette} & \textbf{Calinski-H} & \textbf{Davies-B} & \textbf{Time (s)} \\
\midrule
K-Means (random init) & 0.3465 & 46,712.3 & 1.2891 & 8.2 \\
K-Means++ & 0.3481 & 46,823.7 & 1.2847 & 8.7 \\
\bottomrule
\end{tabular}
\end{table}

\textbf{Observation:} K-Means++ achieves marginally better cluster quality (higher silhouette, lower Davies-Bouldin) with minimal additional computational cost (0.5s). The improvement confirms that smart initialization reduces local optima.

\subsubsection{Cluster Characteristics}

\textbf{K-Means++ Cluster Analysis (K=2):}

\begin{table}[H]
\centering
\caption{Cluster Characteristics (K-Means++, K=2)}
\begin{tabular}{lcccc}
\toprule
\textbf{Cluster} & \textbf{Size} & \textbf{Percentage} & \textbf{Avg. Vehicles} & \textbf{Injury Crashes (\%)} \\
\midrule
Cluster 0 & 118,543 & 56.7\% & 2.03 & 28.4\% \\
Cluster 1 & 90,732 & 43.3\% & 2.18 & 59.6\% \\
\bottomrule
\end{tabular}
\end{table}

\textbf{Interpretation:}
\begin{itemize}
    \item \textbf{Cluster 0}: Lower-severity crashes (28.4\% injury rate), fewer vehicles involved
    \item \textbf{Cluster 1}: Higher-severity crashes (59.6\% injury rate), more vehicles involved
\end{itemize}

\textbf{Observation:} The K=2 clusters strongly correlate with the binary \texttt{crash\_type} target (injury vs. non-injury), suggesting that K-Means essentially rediscovered the target variable rather than finding novel patterns. This confirms that crash severity is the dominant structure in the feature space.

\subsubsection{Silhouette Plot Analysis}

Silhouette plots visualize the quality of cluster assignments for each sample. Observations:

\begin{itemize}
    \item \textbf{Average silhouette score}: 0.3481 (moderate clustering quality)
    \item \textbf{Cluster 0}: Mean silhouette = 0.38, indicating reasonable cohesion
    \item \textbf{Cluster 1}: Mean silhouette = 0.31, indicating weaker cohesion with more scattered points
    \item \textbf{Negative silhouettes}: 18.7\% of samples have negative scores, indicating potential misclassifications or boundary cases
\end{itemize}

\textbf{Observation:} The moderate silhouette score (0.35) and substantial proportion of negative values indicate overlapping clusters with fuzzy boundaries. This reflects the inherent ambiguity in crash severity—many crashes fall on the borderline between injury and non-injury outcomes.

\subsection{DBSCAN Clustering}

\subsubsection{Methodology}

DBSCAN (Density-Based Spatial Clustering of Applications with Noise) identifies clusters as dense regions separated by sparse areas. Unlike K-Means, DBSCAN:
\begin{itemize}
    \item Does not require pre-specifying K
    \item Can find arbitrary-shaped clusters
    \item Labels sparse points as noise/outliers
\end{itemize}

\textbf{Parameters:}
\begin{itemize}
    \item \texttt{eps} ($\epsilon$): Maximum distance between points in a neighborhood
    \item \texttt{min\_samples}: Minimum points required to form a dense region
\end{itemize}

\subsubsection{Parameter Selection}

\textbf{K-Distance Plot:} Computed 5-nearest neighbor distances and sorted them to identify the "elbow" for $\epsilon$ selection. The 95th percentile distance suggested $\epsilon \approx 1.8$.

\textbf{Grid Search:} Tested combinations:
\begin{itemize}
    \item eps: 0.5, 1.0, 1.5, 2.0
    \item min\_samples: 40, 48, 56 (rule of thumb: $\geq 2 \times$ dimensionality = 48)
\end{itemize}

\textbf{Best Parameters:} eps = 0.5, min\_samples = 48 (highest silhouette score among valid clusterings)

\subsubsection{DBSCAN Results}

\begin{table}[H]
\centering
\caption{DBSCAN Clustering Results (eps=0.5, min\_samples=48)}
\begin{tabular}{lcc}
\toprule
\textbf{Metric} & \textbf{Value} & \textbf{Percentage} \\
\midrule
Number of clusters & 4 & --- \\
Noise points & 187,463 & 89.6\% \\
Clustered points & 21,812 & 10.4\% \\
Largest cluster size & 12,847 & 6.1\% \\
Silhouette score (non-noise) & 0.2134 & --- \\
\bottomrule
\end{tabular}
\end{table}

\textbf{Observation:} DBSCAN labeled 89.6\% of crashes as noise, indicating that the vast majority of data points do not fall into dense regions. The 4 small clusters (totaling only 10.4\% of data) represent rare, tightly-grouped crash scenarios. The low silhouette score (0.21) and high noise percentage suggest that traffic crashes do not exhibit strong density-based clustering patterns.

\textbf{Interpretation:} The failure of DBSCAN to find meaningful clusters confirms that crashes are relatively homogeneous in feature space, with no distinct dense regions corresponding to specialized accident types. This is consistent with the moderate granularity of categorical features—without precise spatial or temporal coordinates, most crashes appear similar.

\subsection{Association Rule Mining (Apriori Algorithm)}

\subsubsection{Methodology}

The Apriori algorithm discovers frequent itemsets (feature combinations that occur together) and generates association rules of the form:

\begin{equation}
\text{IF } \{A, B\} \text{ THEN } \{C\}
\end{equation}

\textbf{Key Metrics:}
\begin{itemize}
    \item \textbf{Support}: Proportion of transactions containing the itemset
    \begin{equation}
    \text{Support}(X) = \frac{\text{Count}(X)}{N}
    \end{equation}

    \item \textbf{Confidence}: Conditional probability of consequent given antecedent
    \begin{equation}
    \text{Confidence}(X \rightarrow Y) = \frac{\text{Support}(X \cup Y)}{\text{Support}(X)}
    \end{equation}

    \item \textbf{Lift}: Ratio of observed to expected co-occurrence
    \begin{equation}
    \text{Lift}(X \rightarrow Y) = \frac{\text{Support}(X \cup Y)}{\text{Support}(X) \times \text{Support}(Y)}
    \end{equation}
\end{itemize}

Lift $>$ 1 indicates positive association (co-occurrence is more likely than random chance).

\subsubsection{Transaction Encoding}

Traffic crashes were converted to transactions by encoding each feature-value pair as an item:

\textbf{Example Transaction:}
\begin{verbatim}
[weather_condition=RAIN, lighting_condition=DAYLIGHT,
 first_crash_type=REAR END, crash_type=NO INJURY / DRIVE AWAY, ...]
\end{verbatim}

\textbf{Features Included:} \texttt{weather\_condition}, \texttt{lighting\_condition}, \texttt{first\_crash\_type}, \texttt{trafficway\_type}, \texttt{roadway\_surface\_cond}, \texttt{prim\_contributory\_cause}, \texttt{crash\_type}, \texttt{most\_severe\_injury}

\textbf{Total Transactions:} 209,275 (one per crash)

\subsubsection{Frequent Itemsets}

\textbf{Parameters:}
\begin{itemize}
    \item Minimum support: 5\% (itemset must appear in at least 10,464 crashes)
    \item Maximum itemset length: 3 items
\end{itemize}

\textbf{Results:} 487 frequent itemsets discovered

\textbf{Top 5 Frequent Itemsets:}
\begin{enumerate}
    \item \texttt{weather\_condition=CLEAR}: 78.7\% support
    \item \texttt{crash\_type=NO INJURY}: 43.4\% support
    \item \texttt{lighting\_condition=DAYLIGHT}: 65.2\% support
    \item \texttt{roadway\_surface\_cond=DRY}: 52.8\% support
    \item \texttt{weather=CLEAR, lighting=DAYLIGHT}: 51.3\% support
\end{enumerate}

\textbf{Observation:} The most frequent itemsets correspond to the most common individual feature values, confirming that crashes are concentrated in typical conditions (clear weather, daylight, dry roads).

\subsubsection{Association Rules}

\textbf{Parameters:}
\begin{itemize}
    \item Minimum confidence: 50\% (consequent must occur in at least 50\% of antecedent cases)
\end{itemize}

\textbf{Results:} 1,037 association rules discovered

\textbf{Top 10 Association Rules (by Lift):}

\begin{table}[H]
\centering
\caption{Top 10 Association Rules by Lift}
\scriptsize
\begin{tabular}{lp{4cm}p{4cm}cccc}
\toprule
\textbf{Rule} & \textbf{Antecedent (IF)} & \textbf{Consequent (THEN)} & \textbf{Supp.} & \textbf{Conf.} & \textbf{Lift} \\
\midrule
R1 & Wet Road, No Injury & Rain & 0.050 & 0.638 & 6.16 \\
R2 & Rain & Wet Road & 0.100 & 0.964 & 6.13 \\
R3 & Wet Road & Rain & 0.100 & 0.636 & 6.13 \\
R4 & No Injury Indication, Wet Road & Rain & 0.071 & 0.634 & 6.11 \\
R5 & Rain & No Injury Indication, Wet Road & 0.071 & 0.687 & 6.11 \\
R6 & No Injury Indication, Rain & Wet Road & 0.071 & 0.961 & 6.11 \\
R7 & Daylight, Rain & Wet Road & 0.052 & 0.959 & 6.10 \\
R8 & Rain, No Injury & Wet Road & 0.050 & 0.954 & 6.07 \\
R9 & Daylight, Wet Road & Rain & 0.052 & 0.613 & 5.91 \\
R10 & Rain & Daylight, Wet Road & 0.052 & 0.503 & 5.91 \\
\bottomrule
\end{tabular}
\end{table}

\subsubsection{Rule Interpretation}

\textbf{Rule R2 (Strongest by Confidence):}
\begin{itemize}
    \item \textbf{IF} weather\_condition = RAIN
    \item \textbf{THEN} roadway\_surface\_cond = WET
    \item Support: 9.99\% (occurs in 20,905 crashes)
    \item Confidence: 96.38\% (when it rains, road is wet 96.4\% of the time)
    \item Lift: 6.13 (rain makes wet roads 6.13 times more likely)
\end{itemize}

\textbf{Observation:} This rule is intuitive and validates the Apriori algorithm—rain strongly predicts wet road surfaces with 96\% confidence. The high lift (6.13) confirms that this association is much stronger than random co-occurrence.

\textbf{Rule R12 (Behavioral Pattern):}
\begin{itemize}
    \item \textbf{IF} prim\_contributory\_cause = FOLLOWING TOO CLOSELY
    \item \textbf{THEN} first\_crash\_type = REAR END
    \item Support: 7.85\% (16,438 crashes)
    \item Confidence: 86.18\%
    \item Lift: 4.29
\end{itemize}

\textbf{Observation:} Following too closely leads to rear-end collisions in 86\% of cases, with 4.3x higher likelihood than baseline. This actionable insight suggests that tailgating prevention campaigns could reduce rear-end crashes.

\textbf{Rules Involving Injuries:}
Most high-lift rules involve weather-road surface relationships rather than injury severity. Rules directly predicting injuries have lower lift (2.2-2.8), indicating weaker associations between environmental conditions and injury outcomes.

\textbf{Observation:} The strongest associations are between highly correlated features (rain-wet, following closely-rear end) rather than novel patterns. This suggests limited feature granularity prevents discovery of nuanced injury predictors.

\subsection{Phase IV Summary}

\textbf{K-Means/K-Means++ Clustering:}
\begin{itemize}
    \item Optimal K = 2 clusters (silhouette score: 0.348)
    \item Clusters correspond to low-severity (28\% injury) and high-severity (60\% injury) crashes
    \item K-Means++ provides better initialization than random K-Means
    \item Clustering rediscovered the binary target rather than novel patterns
\end{itemize}

\textbf{DBSCAN Clustering:}
\begin{itemize}
    \item Identified 4 small dense clusters (10.4\% of data)
    \item 89.6\% of crashes labeled as noise
    \item Low silhouette score (0.21) indicates weak density-based structure
    \item Confirms that crashes are relatively homogeneous in feature space
\end{itemize}

\textbf{Association Rule Mining:}
\begin{itemize}
    \item Discovered 487 frequent itemsets and 1,037 association rules
    \item Strongest rule: RAIN $\rightarrow$ WET road surface (confidence: 96.4\%, lift: 6.13)
    \item Behavioral rule: FOLLOWING TOO CLOSELY $\rightarrow$ REAR END collision (confidence: 86.2\%, lift: 4.29)
    \item Most high-lift rules involve obvious correlations (weather-road surface) rather than injury predictors
\end{itemize}

\textbf{Key Insights:}
\begin{enumerate}
    \item Traffic crashes naturally separate into two groups (injury vs. non-injury), validating binary classification approach
    \item Density-based clustering fails due to homogeneity in feature distributions
    \item Association rules confirm domain knowledge (rain causes wet roads, tailgating causes rear-ends)
    \item Limited feature granularity prevents discovery of novel patterns beyond obvious correlations
    \item Unsupervised methods validate supervised learning findings but do not reveal hidden crash types
\end{enumerate}

\section{Recommendations and Conclusions}

\subsection{Best Classifier Recommendation}

Based on comprehensive evaluation across all phases, \textbf{Random Forest with Bagging} is recommended as the optimal classifier for predicting traffic accident injury severity.

\textbf{Justification:}
\begin{enumerate}
    \item \textbf{Highest Accuracy:} 74.06\% on test set, outperforming all other classifiers
    \item \textbf{Balanced Performance:} F1-score of 74.03\% indicates good balance between precision (74.01\%) and recall (74.06\%)
    \item \textbf{Strong AUC:} 0.8101 demonstrates excellent ranking of predictions
    \item \textbf{Robust Cross-Validation:} 74.10\% $\pm$ 0.25\% shows low variance and consistent performance
    \item \textbf{Handles Non-Linearity:} Captures complex feature interactions without explicit feature engineering
    \item \textbf{Interpretability:} Feature importance scores provide actionable insights for policy
    \item \textbf{Scalability:} Parallel tree training enables efficient processing of large datasets
    \item \textbf{No Assumption Violations:} Works with mixed data types, collinear features, and imbalanced classes
\end{enumerate}

\textbf{Optimal Hyperparameters:}
\begin{itemize}
    \item n\_estimators: 200 (ensemble of 200 decision trees)
    \item max\_depth: 15 (limits tree depth to prevent overfitting)
    \item min\_samples\_split: 10 (requires 10 samples to split nodes)
\end{itemize}

\textbf{Deployment Recommendations:}
\begin{itemize}
    \item Deploy as REST API for real-time crash severity prediction
    \item Integrate with dispatch systems to prioritize emergency response
    \item Use predicted probabilities (not just binary predictions) to rank urgency
    \item Retrain quarterly with new crash data to maintain performance
    \item Monitor for data drift (changes in feature distributions over time)
\end{itemize}

\subsection{What Was Learned from This Project}

This project provided hands-on experience with the complete machine learning pipeline:

\textbf{Technical Skills:}
\begin{itemize}
    \item \textbf{Feature Engineering:} Creating domain-driven features (\texttt{is\_rush\_hour}, \texttt{failure\_to\_obey}) significantly improved model performance
    \item \textbf{Dimensionality Reduction:} PCA, SVD, LDA, and VIF analysis revealed that variance is broadly distributed, justifying use of full feature sets
    \item \textbf{Encoding Strategies:} Label encoding, ordinal encoding, and one-hot encoding trade-offs for categorical variables
    \item \textbf{Hyperparameter Tuning:} Grid search optimization improved classifier performance by 3-5\% over default parameters
    \item \textbf{Evaluation Metrics:} Understanding when to prioritize accuracy vs. F1-score vs. AUC based on application requirements
    \item \textbf{Ensemble Methods:} Bagging (Random Forest) and boosting (AdaBoost) consistently outperform single models
    \item \textbf{Clustering Analysis:} K-Means and DBSCAN revealed limited natural clustering in traffic crash data
    \item \textbf{Association Rules:} Apriori algorithm confirmed domain knowledge and identified behavioral patterns
\end{itemize}

\textbf{Domain Insights:}
\begin{itemize}
    \item Crash type, damage level, and primary contributory cause are the strongest predictors of injury severity
    \item Environmental conditions (weather, lighting, road surface) have weaker direct relationships with injuries
    \item Behavioral violations (following too closely, failing to yield) strongly predict specific crash types
    \item Traffic crashes exhibit binary structure (injury vs. non-injury) rather than multiple distinct types
    \item Rain-wet road associations are nearly deterministic (96\% confidence)
\end{itemize}

\textbf{Practical Lessons:}
\begin{itemize}
    \item \textbf{Data quality matters more than algorithm sophistication:} Limited feature granularity constrains all models to 74\% accuracy ceiling
    \item \textbf{Tree-based models excel on tabular data:} Random Forest and AdaBoost outperformed deep learning and SVMs
    \item \textbf{Linear models fail with non-linear boundaries:} LDA, Logistic Regression, and SVM achieved only 52-66\% accuracy
    \item \textbf{Cross-validation is essential:} Single train-test splits can be misleading; 5-fold CV provides robust estimates
    \item \textbf{Interpretability aids deployment:} Feature importance from Random Forest helps stakeholders trust and understand predictions
\end{itemize}

\subsection{Feature Associations with Target Variable}

The most important features associated with \texttt{crash\_type} (injury severity), ranked by Random Forest importance:

\begin{enumerate}
    \item \textbf{first\_crash\_type} (18.47\%): Head-on, angle, and fixed object crashes are high-severity
    \item \textbf{damage} (16.24\%): Higher property damage strongly correlates with injuries
    \item \textbf{prim\_contributory\_cause} (13.89\%): Speeding, reckless driving, and DUI increase injury risk
    \item \textbf{num\_units} (8.92\%): Multi-vehicle crashes have higher injury probability
    \item \textbf{failure\_to\_obey} (6.81\%): Traffic signal/sign violations predict severe crashes
    \item \textbf{trafficway\_type} (5.74\%): Complex intersections (4-way, 5-way) increase severity
    \item \textbf{weather\_condition} (5.23\%): Adverse weather (snow, ice) raises injury risk
    \item \textbf{roadway\_surface\_cond} (4.87\%): Wet, icy roads contribute to severe crashes
    \item \textbf{crash\_hour} (4.56\%): Late-night crashes (10 PM - 4 AM) are more severe
    \item \textbf{lighting\_condition} (4.12\%): Darkness increases injury likelihood
\end{enumerate}

\textbf{Statistical Confirmation:} All 10 features had regression coefficients with p $<$ 0.001 in Phase II, confirming statistical significance.

\subsection{Number of Clusters in Feature Space}

Unsupervised clustering analysis identified:

\textbf{K-Means Optimal K:} 2 clusters
\begin{itemize}
    \item Cluster 0 (56.7\%): Lower-severity crashes (28.4\% injury rate)
    \item Cluster 1 (43.3\%): Higher-severity crashes (59.6\% injury rate)
\end{itemize}

\textbf{Silhouette Score:} 0.348 (moderate cluster quality)

\textbf{DBSCAN Results:} 4 small dense clusters + 89.6\% noise

\textbf{Interpretation:} The feature space exhibits binary structure corresponding to injury vs. non-injury outcomes. No additional crash subtypes (e.g., pedestrian-specific, weather-specific) emerge as distinct clusters, likely due to limited feature granularity. The moderate silhouette score and high DBSCAN noise percentage confirm that crash types form overlapping distributions rather than well-separated groups.

\subsection{Data Granularity Limitations}

A critical finding of this project is that \textbf{insufficient data granularity establishes an effective performance ceiling} for all models. No classifier exceeded 74\% accuracy despite extensive hyperparameter optimization and ensemble methods.

\textbf{Specific Granularity Issues:}

\textbf{1. Binary Target Variable:}
\begin{itemize}
    \item \texttt{crash\_type} reduces complex injury outcomes to just two classes
    \item Real-world severity exists on a spectrum (no injury $\rightarrow$ minor $\rightarrow$ severe $\rightarrow$ fatal)
    \item Binary classification conflates fender benders with catastrophic multi-vehicle crashes
\end{itemize}

\textbf{2. Coarse Temporal Features:}
\begin{itemize}
    \item \texttt{crash\_hour}: Only integer hour (0-23), no minutes or seconds
    \item Missing: Sequential time information (e.g., "5 minutes after previous crash at same location")
    \item Missing: Seasonal trends (day-of-year), holiday effects, daylight saving transitions
\end{itemize}

\textbf{3. Broad Categorical Features:}
\begin{itemize}
    \item \texttt{weather\_condition}: "RAIN" does not distinguish drizzle (minimal risk) from downpour (high risk)
    \item Missing: Temperature, wind speed, visibility distance, precipitation rate
    \item \texttt{lighting\_condition}: "DARKNESS" does not specify streetlight presence/functionality
    \item Missing: Actual illuminance (lux), glare conditions
\end{itemize}

\textbf{4. No Spatial Data:}
\begin{itemize}
    \item Missing: GPS coordinates (latitude, longitude)
    \item Missing: Road geometry (curvature, grade, intersection angle)
    \item Missing: Spatial features (distance to hospital, proximity to high-crash corridors)
    \item Cannot identify specific dangerous intersections or build location-based risk models
\end{itemize}

\textbf{5. Missing Continuous Variables:}
\begin{itemize}
    \item Missing: Vehicle speeds at impact (critical for injury severity)
    \item Missing: Vehicle age, type, safety features (airbags, ABS)
    \item Missing: Driver age, experience, license status
    \item Missing: Traffic volume, congestion level
    \item Missing: Road width, lane count, speed limit
\end{itemize}

\textbf{Impact on Model Performance:}

These limitations explain key findings:
\begin{itemize}
    \item \textbf{Low regression $R^2$ (6.33\%)}: Critical predictors (speed, driver age) are absent
    \item \textbf{Classification ceiling (74\%)}: Many crashes with identical features have different outcomes due to unmeasured variables
    \item \textbf{Weak clustering (silhouette 0.35)}: Coarse features create overlapping distributions
    \item \textbf{Obvious association rules}: Only strong correlations (rain-wet) emerge; nuanced patterns require finer data
\end{itemize}

\textbf{Evidence of Ceiling Effect:}

Across 10 classifiers with extensive hyperparameter tuning, performance clustered tightly:
\begin{itemize}
    \item Best: Random Forest (74.06\%)
    \item Second-best: AdaBoost (74.00\%)
    \item Third-best: Decision Tree (73.31\%)
\end{itemize}

The narrow range (73-74\%) despite vastly different algorithms (tree-based, boosting, neural networks) suggests that all models are extracting the maximum information available from the features. No algorithm can learn relationships that require unmeasured variables.

\subsection{Future Work and Improvement Strategies}

To improve classification performance beyond the current 74\% ceiling:

\textbf{1. Enhanced Data Collection:}
\begin{itemize}
    \item \textbf{GPS coordinates}: Enable spatial modeling, intersection-specific risk scores, distance-to-hospital features
    \item \textbf{Vehicle telemetry}: Impact speed, braking distance, airbag deployment, seatbelt usage
    \item \textbf{Driver demographics}: Age, gender, license history, prior violations, DUI convictions
    \item \textbf{Detailed weather}: Temperature, precipitation rate (mm/hr), wind speed, visibility distance, road friction coefficient
    \item \textbf{Traffic context}: Volume (vehicles/hr), congestion level, time since previous crash at location
    \item \textbf{Fine-grained time}: Timestamp to the second, allowing sequential modeling
\end{itemize}

\textbf{2. Advanced Modeling Techniques:}
\begin{itemize}
    \item \textbf{Deep learning with embeddings:} Neural networks with categorical embeddings for high-cardinality features
    \item \textbf{Gradient Boosting Machines (GBMs):} XGBoost, LightGBM, CatBoost often outperform Random Forest
    \item \textbf{Multi-class classification:} Predict 5-level injury severity (\texttt{most\_severe\_injury}) instead of binary \texttt{crash\_type}
    \item \textbf{Ordinal regression:} Model injury severity as ordered categories rather than independent classes
    \item \textbf{Ensemble stacking:} Combine Random Forest, AdaBoost, and Neural Network predictions with meta-learner
    \item \textbf{Poisson regression:} For count targets like \texttt{injuries\_total}, use Poisson or Negative Binomial regression
\end{itemize}

\textbf{3. Feature Engineering:}
\begin{itemize}
    \item \textbf{Geospatial features:} Distance to nearest hospital, trauma center level, traffic density by zip code
    \item \textbf{Temporal aggregations:} Rolling averages (e.g., "crashes in past 7 days at this location")
    \item \textbf{Interaction features:} Weather $\times$ road surface, hour $\times$ day of week
    \item \textbf{Text mining:} NLP on crash narratives (if available) to extract additional context
\end{itemize}

\textbf{4. Addressing Class Imbalance:}
\begin{itemize}
    \item \textbf{SMOTE:} Synthetic Minority Over-sampling Technique for injury class
    \item \textbf{Class weights:} Penalize misclassification of minority class more heavily
    \item \textbf{Threshold tuning:} Optimize classification threshold for operational objectives (e.g., prioritize recall)
\end{itemize}

\textbf{5. Operational Deployment:}
\begin{itemize}
    \item \textbf{Real-time prediction API:} Integrate with dispatch systems for immediate triage
    \item \textbf{Probabilistic outputs:} Use predicted probability of injury (0-1) instead of binary classification for nuanced decision-making
    \item \textbf{Explainability:} Implement SHAP (SHapley Additive exPlanations) to explain individual predictions to dispatchers
    \item \textbf{Continuous learning:} Retrain models monthly/quarterly as new crash data accumulates
    \item \textbf{A/B testing:} Compare dispatch outcomes with and without model predictions to measure real-world impact
\end{itemize}

\subsection{Final Conclusions}

This comprehensive machine learning project successfully developed and evaluated predictive models for traffic accident injury severity across four analytical phases:

\textbf{Phase I} revealed that crash type, damage level, and primary contributory cause are the most important predictors, accounting for 48.6\% of Random Forest feature importance.

\textbf{Phase II} demonstrated that multiple linear regression explains only 6.33\% of variance in injury counts, indicating weak linear relationships and suggesting that critical predictors are missing from the dataset.

\textbf{Phase III} identified Random Forest with bagging as the best classifier, achieving 74.06\% accuracy, 74.03\% F1-score, and 0.81 AUC, outperforming nine other algorithms including LDA, Decision Trees, Logistic Regression, KNN, SVM, Naive Bayes, stacking, AdaBoost, and Neural Networks.

\textbf{Phase IV} found that traffic crashes naturally cluster into two groups (injury vs. non-injury) with silhouette score of 0.348, and association rule mining discovered 1,037 rules with the strongest being RAIN $\rightarrow$ WET road surface (confidence: 96.4\%, lift: 6.13).

\textbf{Critical Limitation:} Insufficient data granularity—particularly the binary target, coarse temporal features, missing spatial coordinates, and lack of continuous variables (speed, driver age, traffic volume)—establishes a performance ceiling around 74\% accuracy that no algorithm or hyperparameter tuning can overcome.

\textbf{Practical Impact:} Despite limitations, the Random Forest classifier provides actionable predictions for emergency dispatch systems, with 74\% accuracy representing a substantial improvement over random guessing (59\% for majority class baseline). Feature importance analysis guides policy interventions: enforcement of traffic signals (failure\_to\_obey), rear-end collision prevention (following too closely), and targeted safety improvements at complex intersections.

\textbf{Recommendations:} Deploy the Random Forest model as a real-time prediction API for emergency services, prioritizing response to crashes predicted as high-severity. Simultaneously, invest in enhanced data collection (GPS coordinates, vehicle telemetry, driver demographics, detailed weather) to unlock higher prediction accuracy in future iterations.

This project demonstrates the complete machine learning workflow from exploratory data analysis through model deployment recommendations, providing both technical expertise and domain insights for traffic safety applications.

\section{References}

\begin{enumerate}
    \item Breiman, L. (2001). Random Forests. \textit{Machine Learning}, 45(1), 5-32.
    \item Pedregosa, F., et al. (2011). Scikit-learn: Machine Learning in Python. \textit{Journal of Machine Learning Research}, 12, 2825-2830.
    \item Hastie, T., Tibshirani, R., \& Friedman, J. (2009). \textit{The Elements of Statistical Learning: Data Mining, Inference, and Prediction} (2nd ed.). Springer.
    \item National Highway Traffic Safety Administration (2023). \textit{Traffic Safety Facts 2023}. U.S. Department of Transportation.
    \item Agrawal, R., \& Srikant, R. (1994). Fast Algorithms for Mining Association Rules. \textit{Proceedings of the 20th VLDB Conference}, 487-499.
    \item Ester, M., Kriegel, H. P., Sander, J., \& Xu, X. (1996). A Density-Based Algorithm for Discovering Clusters in Large Spatial Databases with Noise. \textit{Proceedings of KDD}, 96(34), 226-231.
    \item Friedman, J. H. (2001). Greedy Function Approximation: A Gradient Boosting Machine. \textit{Annals of Statistics}, 29(5), 1189-1232.
    \item Bishop, C. M. (2006). \textit{Pattern Recognition and Machine Learning}. Springer.
    \item World Health Organization (2023). \textit{Global Status Report on Road Safety 2023}. WHO Press.
    \item Rousseeuw, P. J. (1987). Silhouettes: A Graphical Aid to the Interpretation and Validation of Cluster Analysis. \textit{Journal of Computational and Applied Mathematics}, 20, 53-65.
\end{enumerate}

\section{Appendix: Python Code}

The complete Python implementation for all four phases is provided in the following files:

\begin{itemize}
    \item \texttt{ml\_pipeline\_optimized.py}: Integrated pipeline for Phase I-III
    \item \texttt{phase\_II\_complete.py}: Detailed regression analysis with statistical tests
    \item \texttt{phase\_III\_complete.py}: Comprehensive classification with all 10 algorithms
    \item \texttt{phase\_IV\_complete.py}: Clustering and association rule mining
\end{itemize}

\textbf{Code Repository:} Available upon request. All code follows PEP 8 style guidelines and includes comprehensive documentation.

\textbf{Requirements:}
\begin{verbatim}
pandas==1.5.3
numpy==1.24.2
scikit-learn==1.2.2
matplotlib==3.7.1
seaborn==0.12.2
mlxtend==0.22.0
category-encoders==2.6.0
\end{verbatim}

\textbf{Usage:}
\begin{verbatim}
# Phase I-III (integrated)
python ml_pipeline_optimized.py

# Phase II (detailed regression)
python phase_II_complete.py

# Phase III (all classifiers)
python phase_III_complete.py

# Phase IV (clustering + association rules)
python phase_IV_complete.py
\end{verbatim}

All outputs (figures, tables, metrics) are saved to \texttt{./new\_ver\_results/} directory.

\end{document}
